% Schéma matice proměnných X v převodu NzMna na SAT. Řádky odpovídají pořadí
% prvku množiny, sloupce vrcholům. Barvy podmínek jsou stejné jako v obrázku
% nzmna-sat-formule.tex; podmínky se odkrývají v pořadí, v jakém jsou
% uvedeny na snímku „Podmínky kladené na matici X“ (snímky 3-6). Bubliny
% mají styl formulacallout z assets/macros.tex.
\begin{tikzpicture}[
    pruh/.style={rounded corners=4pt, line width=0.4pt}
]

    % ------------------------------------------------------------------
    % Zvýrazněné řádky a sloupec (pod prvky matice)
    % ------------------------------------------------------------------
    % sloupec j: nejvýše jedna jednička v každém sloupci
    \draw[pruh, visible on=<5->, fill=calloutorangebg, draw=calloutorangefg]
        (3.15,0.38) rectangle (4.35,-4.13);
    % řádek i: nejvýše jedna jednička v každém řádku
    \draw[pruh, fill opacity=0.8, visible on=<3->, fill=calloutredbg, draw=calloutredfg]
        (-0.55,-1.93) rectangle (6.8,-2.57);
    % řádek 1: alespoň jedna jednička v každém řádku
    \draw[pruh, fill opacity=0.8, visible on=<4->, fill=calloutgreenbg, draw=calloutgreenfg]
        (-0.55,0.32) rectangle (6.8,-0.32);

    % ------------------------------------------------------------------
    % Prvky matice: řádky 1, 2, ..., i, ..., k a sloupce 1, 2, ..., j, ..., n
    % ------------------------------------------------------------------
    \foreach \r/\rl/\rd in {1/1/0, 2/2/0, 3/{}/1, 4/i/0, 5/{}/1, 6/k/0} {
        \foreach \c/\cl/\cd in {1/1/0, 2/2/0, 3/{}/1, 4/j/0, 5/{}/1, 6/n/0} {
            \node[visible on=<2->] at ({(\c-1)*1.25}, {(1-\r)*0.75})
                {\ifnum\rd=1
                    \ifnum\cd=1 $\ddots$\else $\vdots$\fi
                 \else
                    \ifnum\cd=1 $\cdots$\else $x_{\rl,\cl}$\fi
                 \fi};
        }
        % číslo řádku = pořadí prvku množiny
        \node[visible on=<2->] at (-1.35, {(1-\r)*0.75})
            {\ifnum\rd=1 $\vdots$\else $\rl.$\fi};
    }

    % závorky matice
    \draw[thick, visible on=<2->] (-0.65,0.45)
        .. controls (-0.95,-0.3) and (-0.95,-3.45) .. (-0.65,-4.2);
    \draw[thick, visible on=<2->] (6.9,0.45)
        .. controls (7.2,-0.3) and (7.2,-3.45) .. (6.9,-4.2);

    % sloupce = vrcholy (proměnné v_1, ..., v_n)
    \foreach \c/\cl/\cd in {1/1/0, 2/2/0, 3/{}/1, 4/j/0, 5/{}/1, 6/n/0} {
        \node[visible on=<2->] (v\c) at ({(\c-1)*1.25}, 1.0)
            {\ifnum\cd=1 $\cdots$\else $v_{\cl}$\fi};
    }
    \node[visible on=<2->, rotate=90, font=\small] at (-2.0,-1.875) {pořadí prvku};

    % ------------------------------------------------------------------
    % Bubliny s podmínkami (texty jako v obrázku výsledné formule)
    % ------------------------------------------------------------------
    \node[formulacallout=red, visible on=<3->,
        callout absolute pointer={(6.8,-2.25)}]
        at (6.8,-2.25) {Nejvýše jedna jednička\\v~každém řádku};

    \node[formulacallout=green, visible on=<4->,
        callout absolute pointer={(6.8,0)}]
        at (6.8,0) {Alespoň jedna jednička\\v~každém řádku};

    \node[formulacallout=orange, visible on=<5->, anchor=north, xshift=0pt,
        callout absolute pointer={(3.75,-4.13)}]
        at (3.75,-4.8) {Nejvýše jedna jednička v~každém sloupci};

    % vazba x_{i,j} => v_j: jednička ve sloupci j vynutí vybrání vrcholu j
    \draw[pruh, visible on=<6->, fill=calloutpurplebg, draw=calloutpurplefg]
        (3.4,0.72) rectangle (4.1,1.28);
    \node[visible on=<6->] at (3.75,1.0) {$v_j$};
    \draw[-latex, thick, visible on=<6->, draw=calloutpurplefg]
        (3.75,0.38) -- (3.75,0.72);
    \node[formulacallout=purple, visible on=<6->,
        callout absolute pointer={(4.1,1.1)}]
        at (4.1,2.0) {Očíslované vrcholy\\jsou vybrané};

\end{tikzpicture}
